\section{Introduction}

Gaussian avatars provide an explicit spatial support that can be animated and rendered efficiently. Wardrobe control is nevertheless difficult when one shared neural decoder must reconstruct localized, high-frequency changes associated with several garments. In our fixed-avatar setting, the observed failure is not a lack of Gaussian support alone: a sufficiently expressive garment-specific teacher can place residuals on that support, while a shared conditional decoder can smooth, saturate, or collapse the residual field.

We therefore study a narrower question: can reference images control a wardrobe of preregistered seen garments by predicting coefficients in an explicit canonical residual basis? Our setting contains one fixed identity and five seen garments. Garment-specific teachers define canonical residual endpoints; centering these five endpoints yields a basis whose complete centered rank is $K=4$. A lightweight reference-only branch predicts basis coefficients, and the target pose and camera are used only by the existing deformation and rendering stages. The construction separates spatial representation from reference conditioning and makes the initial prediction semantics explicit.

The scope is intentionally bounded. The protocol is view-transductive: each forward pass excludes the target image from its references, but the teachers, basis, and predictor are learned using all four preregistered views. O07 is a held-out garment failure, not a success case. We make no claim about garments outside the registered wardrobe, arbitrary clothing, additional identities, held-out cameras, or semantic understanding of garment geometry.

The current paper candidate motivates three evidence questions. First, does continuous coefficient prediction provide value beyond true outfit-ID lookup and reference-to-endpoint classification? Second, is corrected unsaturated coefficient supervision sufficient, and is linear fusion an economical choice relative to a matched complex alternative? Third, are the apparent gains robust to color counterfactuals, extended perceptual and boundary metrics, and spatial-artifact checks? These questions determine the final wording rather than serving as post-hoc analyses.

\resultplaceholder{Order and finalize the contribution list after the main, hard-lookup, M1--M4, and color outcomes are frozen. Candidate contribution components are: the bounded problem formulation; explicit teacher-derived canonical residual bases; reference-only coefficient control with deterministic mean-garment initialization; and an audit-oriented evaluation against lookup and spatial-artifact alternatives.}
